\documentclass[12pt, a4paper]{article}

% === 1. cmap должен быть самым первым! ===
\usepackage{cmap} 

% === Кодировки и язык ===
\usepackage[utf8]{inputenc}
\usepackage[T2A]{fontenc}
\usepackage[english, russian]{babel}

% === 2. ДОБАВЛЯЕМ xcolor для красивых ссылок ===
\usepackage{xcolor} 

% === Математика ===
\usepackage{amsmath, amssymb, amsthm, mathtools}
\usepackage{bm} 

% === Оформление страницы и ссылок ===
\usepackage{geometry}
\geometry{left=2cm, right=1.5cm, top=2cm, bottom=2cm}

% hyperref должен быть одним из последних!
\usepackage{hyperref}
\hypersetup{
    colorlinks=true,
    linkcolor=blue!70!black,
    citecolor=green!50!black,
    urlcolor=blue!70!black,
    pdfauthor={Ваше Имя},
    pdftitle={Конспект семинаров по Анмеху}
}

% === Настройка теорем и определений ===
% Стиль "plain" (курсив внутри, жирный заголовок) - для теорем
\theoremstyle{plain}
\newtheorem{theorem}{Теорема}[section] % Нумерация: Теорема 1.1, Теорема 1.2...
\newtheorem{lemma}[theorem]{Лемма}
\newtheorem{corollary}[theorem]{Следствие}
\newtheorem{proposition}[theorem]{Утверждение}

% Стиль "definition" (прямой шрифт) - для определений
\theoremstyle{definition}
\newtheorem{definition}{Определение}[section]
\newtheorem{example}{Пример}[section]

% Стиль "remark" (курсив, без жирного заголовка) - для замечаний
\theoremstyle{remark}
\newtheorem{remark}{Замечание}[section]

% Перевод слова "Доказательство" (на случай, если babel не подхватит автоматически)
\renewcommand{\proofname}{Доказательство}

% === Красивые заголовки (опционально) ===
\usepackage{titlesec}
\titleformat{\section}{\Large\bfseries}{\thesection}{1em}{}
\titleformat{\subsection}{\large\bfseries}{\thesubsection}{1em}{}

% === Начало документа ===
\begin{document}

% === Титульный блок ===
\title{\textbf{Аналитическая механика}}
\author{Струков Олег \\ \small Б04-404}
\date{}
\maketitle

% === Оглавление ===
\tableofcontents
\newpage

% ==========================================
% === РАЗДЕЛ 1 ===
% ==========================================
\section{Положения равновесия и теоремы, связанные с ними}

\begin{definition}
    \textbf{Положение равновесия} -- положение системы, в котором она может находиться сколь угодно долго, т.е. $\forall t \hookrightarrow  q(t) = q(0)$ при $\dot{q}(t) = 0  $.
\end{definition}


Система находится в положении равновесия $\Leftrightarrow \delta A_{\text{акт}} = 0$ для виртуального перемещения:
$\displaystyle \frac{d\vec{p}}{dt} = \vec{F}_{\text{акт}} + \vec{F}_{\text{пасс}} = m\ddot{\vec{r}} = 0$. Умножаем скалярно на $\delta \vec{r}$; дальшеля идеальных связей $\hookrightarrow \vec{F}_{\text{пасс}} \cdot \delta\vec{r} = 0$ $\Rightarrow$ $\vec{F}_{\text{акт}} \cdot \delta\vec{r} = 0$ $\Rightarrow$. В общем виде для потенц. $Q_i$:
$$\sum_i Q_i \delta q_i = 0 \quad \Rightarrow \quad Q_i = - \left.\frac{\partial \Pi}{\partial q_i}\right|_{q=q^*} = 0$$


\begin{definition}
    Положение равновесия $q^* = 0$ \textbf{устойчиво}, если $\forall \varepsilon > 0 \; \exists \delta > 0: \; \left\| \begin{pmatrix} q(0) \\ \dot{q}(0) \end{pmatrix} \right\| < \delta \Rightarrow \left\| \begin{pmatrix} q(t) \\ \dot{q}(t) \end{pmatrix} \right\| < \varepsilon \; \forall t$.

    Положение равновесия \textbf{устойчиво асимптотически}, если $\lim\limits_{t \to \infty} \left\| \begin{pmatrix} q(t) \\ \dot{q}(t) \end{pmatrix} \right\| = 0$.
\end{definition}



\begin{theorem}[Лагранж-Дирихле]
Если $\Pi$ консервативных сил в окрестности положения равновесия $q^* = 0$ имеет строгий локальный минимум в $q^* = 0$ $\Rightarrow$ положение равновесия устойчиво. Если же имеет максимум хотя бы по одной из координат $q_i$ $\Rightarrow$ неустойчиво.
\end{theorem}

\begin{theorem}[Ляпунов]
Допустим, $\exists$ ОДУ $\dot{x} = f(x)$. Если в окрестности положения равновесия $B_{\varepsilon}(x^*)$ $\exists$ функция $V$: $V(0) = 0$; $V(x) > 0 \; \forall x \in \mathring{B}_{\varepsilon}(x^*)$, то:\\
а) Если $\dot{V}(x) \leqslant 0$ $\Rightarrow$ положение равновесия $x^*$ устойчиво; \\
б) Если $\dot{V}(x) < 0$ $\Rightarrow$ асимптотически устойчиво.
\end{theorem}

\begin{theorem}[Четаев о неустойчивости]
Если в окрестности положения равновесия существует область $\Omega$: $V(x) > 0 \; \forall x \in \Omega$ и $V(x) = 0 \; \forall x \in \partial\Omega$ и $\forall x \in \Omega$ $\hookrightarrow$ $\dot{V}(x) > 0$ $\Rightarrow$ положение равновесия неустойчиво.
\end{theorem}

Далее используем разложение потенциальной энергии тела вблизи положения равновесия по Тейлору:
$$\Pi(q) = \Pi_0 + \sum_i \left.\frac{\partial \Pi}{\partial q_i}\right|_{q^*} q_i + \frac{1}{2} \sum_{i,j} \left.\frac{\partial^2 \Pi}{\partial q_i \partial q_j}\right|_{q^*} q_i q_j + o(|q|^2)$$

$\Pi_0 = 0$ просто потому что мы можем принять любую П за нулевой уровень, $\sum_i \left.\frac{\partial \Pi}{\partial q_i}\right|_{q^*} q_i = 0$ -- следует из теоремы Лагранжа-Дирихле ну и как будто бы и так очевидно.

\begin{theorem}[Ляпунов 1]
Если $\Pi$ в консервативной системе в положении равновесия $q^* = 0$ не имеет минимума, что установлено из рассмотрения членов 2-го порядка, то положение равновесия неустойчиво.
\end{theorem}

\begin{theorem}[Ляпунов 2]
Если $\Pi$ в консервативной системе в положении равновесия $q^* = 0$ имеет максимум, что установлено из рассмотрения членов наименьшего порядка, то положение равновесия неустойчиво.
\end{theorem}

\begin{theorem}[Барбашин-Красовский]
Рассмотрим систему $\dot{x} = f(x)$, $f(0) = 0$; пусть $\exists V(x)$: $V(0) = 0$ и $V(x) > 0 \; \forall x \neq 0$; $\dot{V}(x) \leqslant 0$. Тогда если в $M = \{x: \dot{V}(x) = 0\}$ нет целых траекторий, кроме $x = 0$ $\Rightarrow$ положение равновесия $x = 0$ асимптотически устойчиво.
\end{theorem}

\begin{definition}
    \textbf{Бифуркационная диаграмма} -- график, показывающий зависимость положениий равновесия от некоторого параметра системы.
\end{definition}



\section{Исследование положений равновесия с помощью диффуров}
\subsection{Линейные и почти линейные системы}

Допустим, у нас имеется линейная система $\dot{x} = Ax$, $A \in \mathbb{R}^{n \times n}$; $x, \dot{x} \in \mathbb{R}^n$; либо линейное приближение в окрестности положения равновесия: $f(0) = 0$; $f(x) = \dot{x}$, $x \in \mathbb{R}^n \to \dot{x} = \left.\frac{\partial f}{\partial x}\right|_{x=0} x + g(x)$, $\frac{\partial f}{\partial x} = \begin{pmatrix} \frac{\partial f_1}{\partial x_1} & \cdots & \frac{\partial f_1}{\partial x_n} \\ \vdots & & \vdots \\ \frac{\partial f_n}{\partial x_1} & \cdots & \frac{\partial f_n}{\partial x_n} \end{pmatrix}$ $\Rightarrow$ зависимость устойчивости от собственных чисел $\lambda_i$:

\begin{enumerate}
    \item $\forall i \hookrightarrow  \operatorname{Re}(\lambda_i) < 0 \Rightarrow$ асимптотическая устойчивость;
    \item $\exists i: \operatorname{Re}(\lambda_i) > 0 \Rightarrow$ неустойчивость;
    \item $\exists i: \operatorname{Re}(\lambda_i) = 0 \Rightarrow$ исследуем дальше (например, прямым методом Ляпунова).
\end{enumerate}

\begin{theorem}[Критерий Рауса--Гурвица]
    $\det(A - \lambda E) = P_n(\lambda) = a_0 \lambda^n + \ldots + a_{n-1}\lambda + a_n = 0$;
\[
\Gamma = \begin{pmatrix}
a_1 & a_3 & a_5 & \cdots \\
a_0 & a_2 & a_4 & \cdots \\
0 & a_1 & a_3 & \cdots \\
\cdots&\cdots&\cdots&\cdots\\
0 &0 &\cdots & a_n
\end{pmatrix}
\]
Чтобы для всех $\lambda_i$ выполнялось $\operatorname{Re}(\lambda_i) < 0$:
\begin{enumerate}
    \item Необходимо $\forall i \to a_i > 0$.
    \item Необходимо и достаточно: главные миноры матрицы Гурвица $\Delta_i > 0$.
\end{enumerate}
\end{theorem}

\subsection{Вынужденные колебания}

Допустим, мы получили диффур $A_{n\times n} \ddot{q} + B \dot{q} + Cq= \begin{pmatrix} a \\ 0 \\ ... \\ 0 \end{pmatrix} \sin(\omega t)$. Для устойчивого положения равновесия необходимо, чтобы квадратичные формы А и С были положительно определены, а В -- положительно полуопределена. Решение складывается из общего и частного, и, если В строго положительно определена, то $\lim_{t \to \infty}q_{\text{общ}}(t) = 0$.

Далее ищем решение в виде $q(t) = D_{n\times 1} e^{i\omega t} \Rightarrow ((i\omega)^2A + i\omega B + C) D e^{i \omega t} = \begin{pmatrix} a \\ 0 \\ ... \\ 0 \end{pmatrix} e^{i \omega t}$

$D = ((i\omega)^2A + i\omega B + C)^{-1} \begin{pmatrix} a \\ 0 \\ ... \\ 0 \end{pmatrix} \qquad \Rightarrow \qquad q_k (t)= |d_{k1}| \sin(\omega t + \arg(d_{k1})) $

\subsection{Малые колебания возле положения равновесия}

\[
T = \frac{1}{2}\dot{q}^T A \dot{q}; \quad \Pi = \frac{1}{2}\sum_{i,j} \frac{\partial^2 \Pi}{\partial q_i \partial q_j} q_i q_j + o(|q|^2) = \frac{1}{2}q^T C q \quad \Rightarrow \quad A\ddot{q} + Cq = 0.
\]

\textbf{Уравнение частот:} $\det(C - A\omega^2) = 0 \to (C - A\omega_i^2)u_i = 0$ (находим амплитудные векторы) $\to q = \sum_i c_i u_i \sin(\omega_i t + \alpha_i)$

Свойства решений:
\begin{enumerate}
    \item $\omega_i$ кратное $\Rightarrow$ им соответствуют разные $u_i$.
    \item $\omega_i = 0 \Rightarrow$ соответствует линейное решение $(c_1 t + c_2)u_i$.
    \item $q = U\theta$, $U = [u_1, \ldots, u_n] \to$ подст. в ур. Лагранжа $\to \ddot{\theta} + \Omega\theta = 0$, $\Omega = \operatorname{diag}(\omega_i^2)$.
    \item $u_i$ ортог. и ЛНЗ по метрике $A$ и $C$.
    \item $q = U\theta, \quad U = [u_1, \ldots, u_n] \to \text{в ур-ние Лагранжа} \to \ddot{\theta} + \Omega\theta = 0, \quad \Omega = \operatorname{diag}(\omega_i^2)$
\end{enumerate}





\newpage

\section{Гамильтониан и первые интегралы}

\textbf{Преобразования Лежандра:} $X(x_1, \ldots, x_n)$ порождает $y_i = \frac{\partial X}{\partial x_i}$. Тогда обратное преобразование:
\[
Y = \sum_{i=1}^{n} x_i y_i - X, \quad x_i = \frac{\partial Y}{\partial y_i} \quad \left(\det\frac{\partial^2 X}{\partial x_i^2} \neq 0\right).
\]


В нашем случае лагранжиан $L$ порождает $p_i = \frac{\partial L}{\partial \dot{q}_i}$ $\Rightarrow$ $H(p, q, t) = \sum_i p_i \dot{q}_i - L$%, $\dot{q}_i = \frac{\partial H}{\partial p_i}$

$\frac{d}{dt}\left(\frac{\partial L}{\partial \dot{q}_i}\right) - \frac{\partial L}{\partial q_i} = 0$ $\Rightarrow$ $\dot{p}_i = \frac{\partial L}{\partial q_i}$

$\frac{\partial H}{\partial q_i} = -\frac{\partial L}{\partial q_i} = -\dot{p}_i$ $\Rightarrow$ 
\[
\begin{cases}
\dot{q}_i = \frac{\partial H}{\partial p_i} \\
\dot{p}_i = -\frac{\partial H}{\partial q_i}
\end{cases}, \quad \frac{\partial H}{\partial \alpha} = -\frac{\partial L}{\partial \alpha} \quad \forall \text{ параметра} \alpha.
\]

\begin{definition}
Функция $f: \mathbb{R}^{n+1} \to \mathbb{R}$ называется \textbf{первым интегралом}, если она сохраняется на решениях системы $\dot{x} = g(x, t)$, т.е.
\[
\frac{df}{dt} = \sum_i \frac{\partial f}{\partial x_i} \dot{x}_i + \frac{\partial f}{\partial t} = 0.
\]    
\end{definition}


Для гамильтоновой системы:
\[
\frac{df}{dt} = \sum_i \left(\frac{\partial f}{\partial q_i} \dot{q}_i + \frac{\partial f}{\partial p_i} \dot{p}_i\right) + \frac{\partial f}{\partial t} = 0, \quad 
\begin{cases}
\dot{q}_i = \frac{\partial H}{\partial p_i} \\
\dot{p}_i = -\frac{\partial H}{\partial q_i}
\end{cases}
\]
\[
\Rightarrow \frac{df}{dt} = \sum_i \left(\frac{\partial f}{\partial q_i} \frac{\partial H}{\partial p_i} - \frac{\partial f}{\partial p_i} \frac{\partial H}{\partial q_i}\right) + \frac{\partial f}{\partial t} = \frac{\partial f}{\partial q} \frac{\partial H}{\partial p} - \frac{\partial f}{\partial p} \frac{\partial H}{\partial q} + \frac{\partial f}{\partial t} = \{f, H\} + \frac{\partial f}{\partial t} = 0.
\]

Крит. ПИ: $\frac{df}{dt} = 0$.

\textbf{Скобки Пуассона:}
\[
\{\varphi, \psi\}_{q,p} = \sum_i \left(\frac{\partial \varphi}{\partial q_i} \frac{\partial \psi}{\partial p_i} - \frac{\partial \varphi}{\partial p_i} \frac{\partial \psi}{\partial q_i}\right) = \frac{\partial \varphi}{\partial q} \frac{\partial \psi}{\partial p} - \frac{\partial \varphi}{\partial p} \frac{\partial \psi}{\partial q}.
\]

\textbf{Свойства:}
\begin{enumerate}
    \item $\{\varphi, \psi\} = -\{\psi, \varphi\} = \{-\varphi, \psi\}$; \quad $\{\varphi, \varphi\} = 0$.
    \item Линейность: $\{\alpha_1 \varphi_1 + \alpha_2 \varphi_2, \varphi_3\} = \alpha_1 \{\varphi_1, \varphi_3\} + \alpha_2 \{\varphi_2, \varphi_3\}$.
    \item $\frac{d}{dt} \{\varphi_1, \varphi_2\} = \left\{\frac{d}{dt} \varphi_1, \varphi_2\right\} + \left\{\varphi_1, \frac{d}{dt} \varphi_2\right\}$.
    \item Свойство Якоби: $\{\{\varphi_1, \varphi_2\}, \varphi_3\} + \{\{\varphi_3, \varphi_1\}, \varphi_2\} + \{\{\varphi_2, \varphi_3\}, \varphi_1\} = 0$.
\end{enumerate}

\begin{theorem}[Якоби--Пуассон]
$\varphi_1$ и $\varphi_2$ -- независимые ПИ $\Rightarrow$ $\varphi_3 = \{\varphi_1, \varphi_2\}$ -- тоже ПИ
   
\end{theorem}




\textbf{Известные структуры $H$ при наличии ПИ:}
\begin{enumerate}
    \item $H(q, p) = h$ -- не зависит от $t$.
    \item $H = H(f_1(q_1, p_1), \ldots, f_n(q_n, p_n)) \Rightarrow f_i(q_i, p_i)$ -- ПИ
    \item \textit{Матрёшка:} $H = H(\ldots, f_3(f_2(f_1(q_1, p_1), q_2, p_2), q_3, p_3), \ldots) \Rightarrow f_i(\alpha_{i-1}, q_i, p_i) = \alpha_i$ -- ПИ
    \item $H = \dfrac{f_1(q_1, p_1) \cdots f_n(q_n, p_n)}{g_1(q_1, p_1) \cdots g_n(q_n, p_n)} \Rightarrow f_i(q_i, p_i) - h \cdot g(q_i, p_i) = \alpha_i$ -- ПИ
\end{enumerate}

\begin{theorem}[Лиувилль]
Если для гамильтоновой системы имеются $n$ ПИ в инволюции ($\{\varphi_i, \varphi_j\} = 0$) $\Rightarrow$ система интегрируется в квадратурах.
    
\end{theorem}

\subsection{Вариационный принцип механики}
\begin{definition}
\textbf{Прямой путь} $q_{\text{пр}}(t)$ -- решение уравнения Лагранжа $\dfrac{d}{dt}\left(\dfrac{\partial L}{\partial \dot{q}}\right) - \dfrac{\partial L}{\partial q} = Q^*$.
    
\end{definition}

\textbf{Окольный путь:} $q_{\text{ок}}(t) = q_{\text{пр}}(t) + \delta q_{\text{пр}} = q_{\text{пр}}(t) + \alpha \cdot \eta(t)$, $\|\eta(t)\| < 1$; $\eta(t_0) = \eta(t_f) = 0$.

\textbf{Принцип Гамильтона:} $S = \displaystyle\int_{t_0}^{t_f} L(q, \dot{q}, t)\, dt$ -- действие. $\delta S = 0 \Leftrightarrow q = q_{\text{пр}}$.

\textbf{Следствие:} $\delta S = 0 \Rightarrow \dfrac{d}{dt}\left(\dfrac{\partial L}{\partial \dot{q}_i}\right) - \dfrac{\partial L}{\partial q_i} = Q_i^*$; $q(t_0) = q_0$, $q(t_f) = q_f$ -- краевая задача.

\begin{theorem}[Нётер]
Имеем однопараметрическое преобразование координат и времени: $\tilde{q} = \tilde{q}(q, t, \alpha)$, $\tilde{t} = \tilde{t}(q, t, \alpha)$, причём при $\alpha = 0$ оно тождественно; имеется обратное; $\dfrac{\partial \tilde{q}}{\partial \alpha}$ и $\dfrac{\partial \tilde{t}}{\partial \alpha}$ непрерывны. Если при этом выполняется симметрия $L(t, q, \dot{q})\, dt = \tilde{L}\!\left(\tilde{t}, \tilde{q}, \dfrac{d\tilde{q}}{d\tilde{t}}\right) d\tilde{t} \Rightarrow$ имеем ПИ вида $\Phi = -\xi H + \eta p$, где $\xi = \left.\dfrac{\partial \tilde{t}}{\partial \alpha}\right|_{\alpha=0}$, $\eta = \left.\dfrac{\partial \tilde{q}}{\partial \alpha}\right|_{\alpha=0}$, $p = \dfrac{\partial L}{\partial \dot{q}}$, $\dot{q} = \dfrac{dq}{d\tilde{t}}\left(\dfrac{dt}{d\tilde{t}}\right)^{-1}$. 
\end{theorem}

\textbf{Интегральные инварианты:} Пуанкаре-Картана $J_{\text{П-К}} = \oint p\,\delta q - H\,\delta t$; Пуанкаре $J_{\text{П}} = \oint p\,\delta q$. Если они существуют $ \Rightarrow$ система гамильтонова.

\begin{theorem}[Ли Хуачжун]
Интеграл произвольного вида $I = \oint M(p, q)\,\delta q + N(p, q)\,\delta p$ является интегральным инвариантом  $\Leftrightarrow I = c \cdot J_n$, $c \neq 0$.
    
\end{theorem}

Значение интегрального инварианта не изменяется с течением времени вдоль траектории.

\begin{definition}
\textbf{Фазовый объём:} $I = \int\!\cdots\!\int \delta q_1 \cdots \delta q_n\, \delta p_1 \cdots \delta p_n$.
    
\end{definition}


\begin{theorem}[Лиувилль]
Фазовый объём в гамильтоновых системах сохраняется, для сохранения фазового объёма системы $\dot{x} = f(x)$ необходимо и достаточно $\operatorname{div} f = 0$.
    
\end{theorem}

\subsection{Канонические преобразования}
\begin{definition}
Канонические преобразования переводят любую гамильтонову систему в гамильтонову систему.
    
\end{definition}

\textbf{Критерий каноничности:} $\oint p\,\delta q - H\,\delta t = c \oint \tilde{p}\,\delta\tilde{q} - \tilde{H}\,\delta t$ $\Rightarrow$ $p\,\delta q - H\,\delta t = c\tilde{p}\,\delta\tilde{q} - c\tilde{H}\,\delta t + \delta F$.

\textbf{Локальный критерий:} $z = \begin{pmatrix} q \\ p \end{pmatrix}$, $J = \begin{pmatrix} 0 & E \\ -E & 0 \end{pmatrix}_{2n \times 2n}$ $\Rightarrow$ $\dot{z} = J \dfrac{\partial H}{\partial z}$ $\Rightarrow$ требуем $\dot{\tilde{z}} = cJ \dfrac{\partial \tilde{H}}{\partial \tilde{z}}$.

\textbf{Важные виды $F$:}
\begin{enumerate}
    \item $F(q, \tilde{q}, t)$ -- свободное.
    \item $F(p, \tilde{p}, t)$
    \item $F(q, \tilde{p})$ 
    \item $F(\tilde{q}, p)$.
\end{enumerate}

\textbf{Эквивалентные соотношения для свободного преобразования:} $c\tilde{p} = -\dfrac{\partial F}{\partial \tilde{q}}$; $p = \dfrac{\partial F}{\partial q}$, $c\tilde{H} = \dfrac{\partial F}{\partial t} + H$.

\textbf{Особый случай:} $\tilde{H} \equiv 0$ $\Rightarrow$ ПИ: $\tilde{q}_i = \alpha_i$, $\tilde{p}_i = \beta_i$.

Возьмём $S(q, \tilde{q}, t)$, запишем для неё критерий каноничности: $p\,\delta q - H\,\delta t = c\tilde{p}\,\delta\tilde{q} + \dfrac{\partial S}{\partial q}\,\delta q + \dfrac{\partial S}{\partial \tilde{q}}\,\delta\tilde{q} + \dfrac{\partial S}{\partial t}\,\delta t$. Тогда $p_i = \dfrac{\partial S}{\partial q_i}$, $\tilde{p}_i = -\dfrac{\partial S}{\partial \tilde{q}_i}$

Получаем \textbf{уравнение Гамильтона-Якоби:} $H + \dfrac{\partial S}{\partial t} = 0$; 

\textbf{Уравнения движения:} $p_i = \dfrac{\partial S}{\partial q_i}$; $\beta_i = \dfrac{\partial S}{\partial \alpha_i}$.






\end{document}